\documentclass[12pt]{article}
\usepackage[T1]{fontenc}
\usepackage[utf8]{inputenc}
\usepackage{lmodern}
\usepackage{textcomp}
\usepackage{enumitem}
\usepackage{geometry}
\geometry{margin=1in}
\begin{document}
\begin{center}
Answer Key - Constitutional Law - Graduate Level Assessment 8 \
\small (Answers are ordered exactly as in the question paper.)
\end{center}

\section*{Multiple Choice}
\begin{enumerate}[label=\arabic*.]
\item \textbf{Answer:} B
\\ \textit{Options:} A) Death or insanity of a party.; B) Partial performance; C) Destruction of the subject matter of contract.; D) Supervening illegality
\\ \textit{Note:} Partial performance does not terminate a contract by operation of law because it indicates that the parties are still fulfilling their obligations under the contract, rather than indicating an end to the contractual relationship.
\item \textbf{Answer:} A
\\ \textit{Options:} A) Warranty of Fitness for Particular Purpose.; B) Warranty of Merchantability; C) Express; D) Implied
\\ \textit{Note:} This is the correct answer because a Warranty of Fitness for Particular Purpose arises when a buyer relies on the seller's expertise to select a product that meets specific needs, indicating that the bicycle must fulfill the buyer's particular goals as understood by the salesperson.
\item \textbf{Answer:} C
\\ \textit{Options:} A) A breaking and entry; B) Of a dwelling, of another; C) During the daytime; D) With the intent to commit a felony therein.
\\ \textit{Note:} The correct answer is "During the daytime" because common law burglary specifically requires the act to be committed at night, distinguishing it from other forms of theft or trespass.
\item \textbf{Answer:} A
\\ \textit{Options:} A) Gender; B) Interstate travel; C) Privacy; D) Alienage
\\ \textit{Note:} The correct answer is gender because regulations based on gender are subject to intermediate scrutiny rather than strict scrutiny, which applies only to classifications based on race or fundamental rights.
\item \textbf{Answer:} B
\\ \textit{Options:} A) No duty for artificial conditions.; B) A duty for natural conditions.; C) No duty for natural conditions.; D) No duty for active operations.
\\ \textit{Note:} Under common law principles, possessors of land do not owe a duty to undiscovered trespassers for natural conditions on the land, as they are not considered foreseeable users of the property.
\end{enumerate}

\section*{True / False}
\begin{enumerate}[label=\arabic*.]
\setcounter{enumi}{5}
\item \textbf{Answer:} True
\\ \textit{Options:} A) True; B) False
\\ \textit{Note:} The Sixth Amendment right to counsel is only applicable after formal charges have been initiated, and investigative surveillance occurs before any charges are filed, meaning an indigent person does not have the right to counsel during that stage.
\item \textbf{Answer:} False
\\ \textit{Options:} A) True; B) False
\\ \textit{Note:} A valid treaty does not always take precedence over congressional acts because the Supremacy Clause of the Constitution establishes that federal law, including laws enacted by Congress, can take precedence over treaties if they conflict, particularly when the law is enacted after the treaty.
\item \textbf{Answer:} False
\\ \textit{Options:} A) True; B) False
\\ \textit{Note:} The enabling clause of the Fourteenth Amendment primarily empowers Congress to enforce the provisions of the amendment against state actions, rather than granting states the authority to regulate private individuals' discrimination based on nationality.
\item \textbf{Answer:} False
\\ \textit{Options:} A) True; B) False
\\ \textit{Note:} Differential treatment of individuals regarding fundamental rights is primarily a matter of equal protection under the law, rather than a substantive due process issue, as substantive due process focuses on the protection of certain fundamental rights from government interference, while equal protection addresses discrimination and unequal treatment.
\item \textbf{Answer:} False
\\ \textit{Options:} A) True; B) False
\\ \textit{Note:} An irrevocable offer under the Uniform Commercial Code must specify a reasonable time period for which it remains open, or it will be deemed invalid if no time frame is established.
\end{enumerate}

\section*{Long Form Response}
\begin{enumerate}[label=\arabic*.]
\setcounter{enumi}{10}
\item \textbf{Answer:} Contract termination by operation of law refers to the automatic cessation of contractual obligations due to specific legal principles or circumstances, such as impossibility, illegality, or mutual agreement. In constitutional law, this principle emphasizes that certain conditions, rather than mere actions taken by the parties, dictate the validity of termination. Partial performance, while demonstrating an effort to fulfill contractual obligations, does not constitute a valid termination method because it does not meet the threshold required for legal termination; it merely indicates that the contract is still in effect, albeit partially executed. Therefore, the legal framework necessitates clear criteria for termination, which partial performance fails to satisfy, reinforcing the need for unequivocal causes for contract dissolution under the law.
\\ \textit{Note:} correct because it emphasizes that contract termination by operation of law relies on specific legal circumstances rather than the actions of the parties, and partial performance does not meet the criteria for valid termination as it does not negate the legal obligations established by the contract.
\item \textbf{Answer:} In common law, the legal principles surrounding attempted crimes hinge on two key elements: intent and substantial steps toward the commission of the crime. In the scenario provided, the defendant demonstrated a clear intention to commit arson by planning to set her home ablaze for insurance proceeds. Although her actions-cooking bacon and leaving the flame on-indicate an effort to execute her plan, the fact that she left the window open, which ultimately prevented the fire from igniting, complicates her claim of an attempt. To establish attempted arson, courts typically require that a defendant not only intends to commit the crime but also takes substantial steps towards its completion, which, in this case, the defendant did. However, because the fire did not occur, the prosecution must prove that her actions were sufficiently proximate to the commission of arson, thus raising questions about whether mere preparation suffices for an attempt under the common law standard.
\\ \textit{Note:} The answer correctly identifies that in common law, for a crime to be considered an attempt, the defendant must exhibit a specific intent to commit the crime and take substantial steps toward its completion, which is illustrated by the defendant's planning and actions even though the crime was not successfully executed.
\end{enumerate}

\vfill
\noindent\textit{End of Answer Key}
\end{document}